\section{Conclusão}

Este trabalho projetou, implementou e validou parcialmente uma bancada de
aquisição para conjuntos motopropulsores com motor BLDC. O sistema é composto
por dois nós embarcados dedicados --- SIBA, para grandezas aerodinâmicas, e
SAME, para grandezas elétricas do motor --- integrados por protocolo UDP a um
supervisório em Python e a um observador por modos deslizantes para estimação
de velocidade, ângulo de rotor e torque. A seguir, cada um dos cinco objetivos
específicos estabelecidos na Introdução é retomado à luz do que foi executado.

\begin{enumerate}

    \item \textbf{Sistema de medição elétrica do motor BLDC (SAME):}
    a placa foi projetada e fabricada com sensores de corrente por núcleo
    toroidal (SS49E), condicionamento de tensão por divisor resistivo e
    condicionamento de corrente por seguidor de tensão. O sensor de corrente
    foi caracterizado em bancada, com linearidade confirmada até
    \SI{516}{\gauss} e sensibilidade dentro da faixa de catálogo. A escala de
    cada canal está definida e o zero do caminho de corrente é levantado a
    cada energização. A cadeia de corrente foi exercitada com o motor real em
    giro, fornecendo as formas de onda de tensão e corrente de fase que
    alimentaram o observador. Permanecem pendentes a verificação da cadeia de
    tensão contra instrumento de referência e a caracterização do piso de
    ruído com o ESC chaveando. Objetivo \textbf{majoritariamente atingido}.

    \item \textbf{Modelagem do motor e projeto do observador:}
    resistência e indutância de fase foram medidas experimentalmente
    (\autoref{tab:resistencia_motor}, \autoref{tab:indutancia_motor}). O
    observador por modos deslizantes com PLL foi projetado, ajustado por
    \textit{autotuning} (evolução diferencial sobre \num{1278} avaliações) e
    validado em simulação. Em bancada, os coeficientes de perdas mecânicas
    foram identificados por regressão sobre \num{19} patamares de velocidade:
    $B = (32{,}4 \pm 4{,}4)\times 10^{-6}$~\si{\newton\metre\second\per\radian}
    e $\tau_c = (15{,}5 \pm 2{,}4)$~\si{\milli\newton\metre}, com
    $R^2 = 0{,}925$. O observador foi validado sobre \num{378} registros reais
    de osciloscópio: erro médio de \SI{0,092}{\percent} em velocidade, desvio
    de \SI{0,3}{\degree} em ângulo e \SI{-7,2}{\percent} em torque quando a
    corrente medida é projetada no ângulo estimado. Objetivo \textbf{atingido}.

    \item \textbf{Sistema de medição aerodinâmica da hélice (SIBA):}
    o condicionamento das células de carga e do tubo de Pitot foi integrado ao
    \textit{firmware} da SIBA com transmissão em tempo real pelo servidor UDP
    embarcado. O tubo de Pitot foi calibrado em túnel de vento pela Equipe
    Albatroz no IPT, com as constantes do sensor embarcadas na SIBA
    (\autoref{sec:calibracao_pitot}, Anexo~\ref{anexo:ensaio_pitot}). O
    procedimento de identificação de $C_T$ e $C_Q$ em bancada está definido
    (\autoref{sec:coef_aero_bancada}); a execução física com hélice acoplada
    permanece pendente. Em simulação, a identificação do coeficiente
    aerodinâmico global $k_Q$ a partir da estimativa de torque do observador
    atingiu erro relativo de \num{3,46}\% frente ao modelo da planta.
    Objetivo \textbf{parcialmente atingido}.

    \item \textbf{Rede de comunicação entre os nós e o supervisório:}
    o protocolo UDP estruturado, a estrutura de comando e o pacote de
    telemetria foram implementados, padronizados por geração de código a
    partir de arquivos YAML comuns ao \textit{firmware} e à aplicação
    \textit{desktop}, e validados de ponta a ponta. Objetivo \textbf{atingido}.

    \item \textbf{\textit{Software} de supervisão em Python:}
    a aplicação cobre monitoramento em tempo real, parametrização interativa,
    calibração de sensores, registro em CSV e diagnóstico via
    \textit{TinyTerminal}. Objetivo \textbf{atingido}.

\end{enumerate}

Em síntese, a infraestrutura de \textit{hardware}, \textit{firmware},
comunicação e \textit{software} de supervisão está integralmente implementada.
A cadeia elétrica --- do terminal do motor até o observador --- foi percorrida
com o motor real e os resultados são compatíveis com o modelo adotado no
projeto. O que distingue os objetivos atingidos do parcialmente atingido é a
ausência do ensaio aerodinâmico completo: hélice acoplada, células de carga e
tubo de Pitot operando simultaneamente sob supervisão da SIBA. Esse ensaio,
que fecha o ciclo da bancada, está planejado para o segundo semestre de
\num{2026} quando a estrutura da equipe estiver disponível.
